
  % How generation interval inference is done in practice

  % There are essentially two tiers of methods:

  % 1. Population-level / naive pooling (most common). All observed generation intervals (or serial
  % intervals) are treated as iid draws from a parametric distribution, typically Gamma. Parameters are
  %  estimated by MLE. This is the "population-level method" in Park et al. 2020 (Section 7.4):
  % estimate the mean and shape of a Gamma from the aggregated data, optionally correcting for
  % right-truncation by exponential reweighting (their Eq. 7.7). This approach is used by the majority
  % of empirical studies — for example, the widely-cited COVID-19 generation interval estimates from
  % Ganyani et al. and others.

  % 2. Individual-level Poisson process models (more sophisticated). Each infector j's transmission is
  % modeled as a non-homogeneous Poisson process with rate proportional to the GI density. This is Park
  %  et al.'s Eq. 7.8–7.11. The likelihood for infector j with n_j offspring at generation intervals
  % s_{i,j} is:

  % L_j ∝ Λ^{n_j} exp(−Λ F₀(t_trunc − t_j)) × ∏ᵢ f₀(s_{i,j})

  % This accounts for right-truncation and the number of offspring. Hart et al. (2022, eLife) extend
  % this further using data-augmentation MCMC for household data, jointly estimating infection times
  % and transmission trees.

  % The critical point: in both approaches, generation intervals from the same infector are treated as
  % conditionally independent. Park et al. acknowledge this explicitly after their Eq. 7.11:

  % "In theory, the forward generation intervals arising from the same infector may be correlated
  % because of non-independence in the contact process [35]; although we do not account for this
  % potential correlation in our likelihood..."

  % They note this is a known gap but find empirically that the iid approximation works adequately for
  % their purposes (estimating R₀ from r).

  % No published method I found accounts for within-infector correlation of the kind your model
  % introduces. The correlation Park et al. reference (from the contact process) is a network effect.
  % Your model introduces a more fundamental source: the shared latent component l_i from the
  % infectiousness profile itself.

  %   Sources:
  % - https://royalsocietypublishing.org/doi/10.1098/rsif.2019.0719
  % - https://elifesciences.org/articles/70767
  % - https://royalsocietypublishing.org/doi/10.1098/rspb.2015.2026
  % - https://royalsocietypublishing.org/doi/10.1098/rsif.2020.0756
  % - https://journals.plos.org/ploscompbiol/article?id=10.1371%2Fjournal.pcbi.1014239
  % - https://pmc.ncbi.nlm.nih.gov/articles/PMC10990235/






















  \clearpage
\section{Discussion}

\begin{itemize}
  \item Brief summary of what we did 
  \item High level: the share of within- vs. between-individual variation in infection timing has a substantial impact on both uncontrolled and controlled epidemics. 
  \item To illustrate this, we introduce a novel modeling framework, the gamma time-shifted model. 
  \item We combine this with a general contact process model that extends to the individual level the basic renewal equation epidemic model. 
  \item We expect this will have implications for gathering contact tracing data and for building projection and forecasting models, highlighting a key point of underexplored uncertainty and setting forth a way of dealing with it parsimoniously, even if just in sensitivity analyses. 
\end{itemize}

\begin{itemize}
  \item Why have we gone so long without really caring about $\psi$? One: it doesn't impact mean-field dynamics, by design. We can get good assessments of the overall behavior of an epidemic without it. Also, we've dealt with it implicitly a lot, we just haven't isolated it as an independent axis of variation
  \item For example: individual-based models make assumptions about $\psi$ all the time, as  do the papers we referenced in the second-to-last paragraph of the intro. 
  \item We can also adjust $\psi$ using methods like the linear chain trick, stringing together sets of dummy exposed and infectious categories; but this also changes the generation interval distribution. Figuring out how to vary the individual ifnectiousness profile without affecting the population-level generation interval distribution is the main contribution of this work. 
\end{itemize}

\begin{itemize}
  \item For the time shift story: the shift in $E[\varsigma]$ is minor; the inflation of $\text{Var}[\varsigma]$ is the real story. 
  \item For the inference of generation intervals and r: this suggets we need to develop good statistical tools and we need to be careful as we collect our data. 
  \item for the overdispersion bit: this a fundamentally new mechanism of superspreading that we should be aware of. A key point: ``super-shedders' and ``super-contacters'' are theoretically identifiable as individuals, but when superspreading arises due to poorly timed infectiousness, that's not necessarily something we can anticipate. It suggests there might be a limit to which superspreading can be targeted, and opens up the need for us to bettter understand the relative contributions of these mechanisms to overdispersion. 
  \end{itemize

\begin{itemize}
  \item  For the sensitivity of TE: note that shifts in symptom timing are commonplace (happend with SARS-CoV-2), and small changes in testing frequency are easy operational decisions. So, understanding the shape of the individual infectiousness profile is important for anticipating how these small shifts might impact epidemic behavior. Also, something we didn't explore here but could be fleshed out more is that we might even use these kinds of shifts, or maybe testing, to assess the shape of the individual infectiousness profile -- they might give us information beyond just simple contact tracing data. 
  \item More details on the counter-intuitive impact of TE/GE contraction/overdispersion for detect-and-isolate, and Reff/overdispersion for gathering size restrictions. We often think of interventions in terms of how they reduce $R_{eff}$, but that's a bit dangerous, because the full effect can work in multiple directions (reduction in $R_{eff}$ but not a commensurate drop in epidemic growth rate, for example). The framework here gives a robust way of dealing with this beyond just naive reductions in R. 
\end{itemize}

\begin{itemize}
  \item Identifiability is a challenge, but seeems theoretically possible under reasonably good circumstances. 
  \item Available data suggests that there may be a spectrum of $\psi$-values across pathogens, though there are huge limitations to the data and the inference methods we've used here, which we discuss in a moment. 
  \item That said, there are reasons to believe there could be differences in $\psi$ across pathogens; some we know generate long-term shedding while others might be short and fast, and our findings roughly align with this taxonomy. 
  \item There are local-level impacts that could influence the observed $\psi$, like local susceptible depletion: for measles, maybe we have a short infectiousness window because everybody gets infected immediately. In theory, a person could go on infecting, but they've already burned through  all their contacts as soon as they reach the infectious stage. That's one thing that could happen and is hard to disentangle with existing data. (mabye epidemiological challenge studies, like when they stuck hamsters in the air vents of people with tuberculosis?) 
\end{itemize}

\begin{itemize} 
  \item Limitations abound: the gamma time shift model is just one choice 
  \item the contact process model is hard to ground empirically 
  \item Lots more could be done with test-based control assessments 
  \item inference is wayyyyy over-simplified; we should talk in detail about all the various confounders that could imapct this, many of which are in the comments at the end of code/psi_empirical_litgi.R. 
\end{itemize }

\begin{itemize}
  \item But, we still think this will have an impact. It opens up new avenues of research and highlights a fundamental epidemiological quantitiy worth studying in its own right. 
  \item We need better data, and better ways of understanding how to use existing data to get at $\psi$. 
  \item We expect this will help generate better projections and forecasts -- or at least projections and forecasts that are more honest about unertainty, and clearar about a potential mechanism driving some of that uncertainty. Thus, we hope it will impact public health practice overall. 
\end{itemize}



























  ---
  1. Benchmark pathogen parameters
  - Table of $R_0$, $\bar{g}$, $\text{Var}[g]$, $\alpha$, $\beta$ for influenza, SARS-CoV-2 omicron,
  and measles. (Referenced as "Table XX" in results.tex line 62 and model.tex line 93; already
  appears in supplementary.tex but needs a main-text version or cross-reference.)
  - Literature sources for each pathogen's parameters.

  2. Simulation details
  - Already summarized in model.tex (simulation approach subsection), but results.tex refers to
  "Methods" for additional details in several places. The current model.tex section points to
  "Supplementary Methods" for full details. Methods should at minimum cross-reference or briefly
  restate the key simulation parameters: $N = 10{,}000$, $n_\text{sim} = 5{,}000$, epidemic
  establishment threshold of 500 infections.

  3. Growth rate estimation procedure
  - Results.tex line 80: "we fit a Poisson generalized linear model to the daily infection counts."
  Methods should specify: infinite population, 7 days of data after 100 cumulative cases, daily
  binning, Poisson GLM with log link and day as covariate.

  4. Generation interval estimation procedure
  - Results.tex lines 82–93: Simulated offspring from index cases, applied Gamma observation noise
  ($a_\text{obs}$, $b_\text{obs}$), extracted serial intervals, estimated posterior for $(\alpha,
  \beta)$ treating observations as independent. Methods should specify: number of index cases (100 in
   results.tex line 93, but 5,000 in uncontrolled.tex line 80), the posterior estimation approach
  (what prior, what likelihood), the coverage calculation (95% credible interval containing the true
  value), and the observation delay parameters ($a_\text{obs} = 4$, $b_\text{obs} = 1$ from
  results.tex line 83).

  5. ICC and effective sample size (Kish's design effect)
  - Results.tex lines 82–89: The ICC formula and $n_\text{eff} = n / (1 + (R_0 - 1) \cdot
  \text{ICC})$. Methods (or supplementary methods) should derive this. The ICC formula itself is
  stated in results.tex but the derivation is promised as "Supplementary Materials."

  6. Overdispersion estimation from simulations
  - Results.tex line 121: "We estimated OD by moment-matching the mean and variance of the offspring
  counts to a Negative Binomial distribution." Methods should state this explicitly.
  - Also specify: 10,000 index cases per scenario, range of $c_\text{amp}$, $\sigma$, $\psi$ values
  swept.

  7. Testing effectiveness (TE) numerical integration
  - Results.tex line 141: "TE can then be calculated via numerical integration (Methods)." This is
  explicitly promised.
  - Controlled.tex lines 19–21 give the integral: $\text{TE} = \eta \rho \int_{-\infty}^\infty (1 -
  F_\epsilon(t)) f_\text{det}(t) , dt$. Methods should specify:
    - The two detection mechanisms (symptom-based and test-based)
    - Symptom-based: $D \sim \text{Normal}(\delta_\text{symp}, \sigma_\text{symp}^2)$ relative to
  peak infectiousness, immediate isolation
    - Test-based: tests at intervals $\Delta_\text{test}$, detectability window $[\delta_\text{pre},
  \delta_\text{post}]$ around peak, $\text{Exp}(\lambda_\text{act})$ turnaround delay
    - How $f_\text{det}$ is constructed for each mechanism
    - Default parameter values used in figures

  8. Detect-and-isolate simulation details
  - Results.tex line 167: Simulation parameters for D&I epidemic simulations ($\rho = 0.5$?, $\eta =
  1$?, $\delta_\text{symp}$, $\sigma_\text{symp}$, $\Delta_\text{test}$, etc.). The controlled.tex
  draft (line 42) specifies $\rho = 0.5$, $\eta = 1$, deterministic isolation 1 day before peak.
  Methods should state these and the comparison setup (naive $R_\text{eff}$-matched simulations).

  9. Generation interval truncation formula
  - Results.tex lines 151–159: $g^*(\tau)$ formula. The derivation is promised as "Supplementary
  Methods."

  10. $\nu_i^*$ under D&I and overdispersion
  - Results.tex lines 161–165: $\nu_i^* = \nu_i F_\epsilon(m_\epsilon + \iota_i)$. Methods or
  supplementary should derive this and the associated dispersion index expressions.

  11. Gathering size restriction details
  - Results.tex lines 174–188: The truncated stochastic contact model, $R_\text{eff} = \mu_T R_0$, OD
   formulas. Methods should specify the truncation mechanism, the values of $\sigma$ and
  $c_\text{max}$ used, and the three-way comparison setup (truncated vs. $R_\text{eff}$-adjusted
  stochastic vs. $R_\text{eff}$-adjusted uniform).

  12. Psi identifiability: simulation-based power calculation
  - Results.tex lines 196–198: 10/25/50/100 index cases, $\psi \in {0, 0.5, 1}$, three pathogens,
  Gamma observation delays (fast/medium/slow), flat prior on $\psi$, 200 replicates, power defined as
   >80% of simulations with 95% posterior mass within 0.2 of true $\psi$. Methods should detail:
    - The posterior computation (likelihood, prior, grid)
    - The power criterion
    - The observation delay parameters (mean 1, 2, 4 days)

  13. Empirical $\psi$ estimation
  - Results.tex lines 209–211 (your new section): Methods should describe:
    - Data extraction from OutbreakTrees database
    - Supplementary COVID-19 data sources (Ganyani, Hart)
    - Cluster formation (grouping by index case, restricting to $\geq 2$ secondaries)
    - Likelihood: integration over shared $l_i$, product of SI densities within cluster
    - Prior: flat on $\psi \in [0, 1]$, grid of 101 points
    - Fixed GI parameters from literature (and their sources, per pathogen)
    - Incubation period parameters used
    - Summary statistics: posterior mode, mean, 95% CI

  14. Items from the older draft that the results reference as "Supplementary Methods"
  These are derivations already partly present in supplementary.tex but should be cross-referenced
  from Methods:
  - Deriving population-level quantities from individual-level components (supplementary.tex line 4)
  - $\text{Var}[W]$ derivation and the Gamma approximation for $W|\text{surv}$ (already in
  supplementary.tex)
  - Overdispersion in daily infection counts (already in supplementary.tex)
  - OD formula for periodic contacts (already in supplementary.tex)
  - OD formula for stochastic contacts (already in supplementary.tex)
  - TE integration details
  - $g^*(\tau)$ derivation
  - Gathering size restriction OD formulas

  15. Software and reproducibility
  - R version, key packages (odin for simulations, tidyverse, etc.)
  - Code availability statement